\section{Binary System}

\subsection{The $X$--$P$--$T$ Diagram}

\paragraph{Composition variable and phase rule.}
A \textcolor{blue}{binary system} has two components, $C=2$, so one
independent variable, $X=x_2$ with $x_1=1-X$, describes the composition. The
phase diagram is therefore three-dimensional in $(P,T,X)$. The diagram can be
pictured as a stack of fixed-composition $P$--$T$ diagrams, one for each
value of $X$. The phase rule reads $F=4-\mathcal{P}$, so each unary
structure gains one dimension:
\begin{itemize}
    \item \textcolor{blue}{Tervariant volume} ($\mathcal{P}=1$, $F=3$). The
    bivariant regions of the unary diagram become volumes in which $P$, $T$
    and $X$ all vary freely without a phase change. The cross section of
    each volume changes with $X$.
    \item \textcolor{blue}{Bivariant surface} ($\mathcal{P}=2$, $F=2$). The
    univariant curves become surfaces. Two coexisting phases generally have
    different compositions, so each phase has its own bivariant surface, and
    tie lines join the two surfaces. The region between the surfaces is the
    two-phase region, in which the overall composition fixes the amount of
    each phase.
    \item \textcolor{blue}{Univariant curve} ($\mathcal{P}=3$, $F=1$). The
    invariant points become curves, along which only one variable varies
    freely. Three bivariant surfaces, one for each pair of phases, meet
    along each such curve.
    \item \textcolor{blue}{Invariant point} ($\mathcal{P}=4$, $F=0$). A new
    kind of invariant point appears, at which none of $P$, $T$ or $X$ varies.
    Every subset of the four phases coexists near this point, so
    $\binom{4}{3}=4$ univariant curves, $\binom{4}{2}=6$ two-phase
    (bivariant) surfaces and $\binom{4}{1}=4$ single-phase (tervariant)
    volumes meet there. This point is therefore also called the
    \textcolor{blue}{quadruple point}.
\end{itemize}
